\chapter{Exponential and Logarithm}

The goal of this chapter is not to borrow exponential or logarithm from an
ambient real analysis library.  It is to construct the algorithms and then
prove that the standard constructive definitions agree.  The three
independent routes are:
\[
  x\longmapsto\sum_{n=0}^{\infty}\frac{x^n}{n!},
  \qquad
  x\longmapsto\lim_{m\to\infty}\left(1+\frac{x}{m}\right)^m,
\]
and the inverse of the logarithmic integral
\[
  \operatorname{Log}(x)=\int_1^x \frac{\dd t}{t}.
\]
The formalization is meant to expose the executable raw functions first.  The
agreement theorems are then organized as finite interval estimates,
calculus uniqueness, and inverse-function certificates.

\section{Power-Series Exponential}

\begin{definition}[Real power-series exponential]\label{def:exp-power-series-real}
\lean{ComputableAnalysis.expPowerSeries}
\uses{def:raw-real}
\leanok
For \(x\in\Q\), define the power-series exponential algorithm by rational
interval stages around the finite partial sums
\[
  S_N(x)=\sum_{k=0}^{N-1}\frac{x^k}{k!}.
\]
The omitted tail is bounded by a geometric majorant.  At stage \(n\), the
algorithm chooses a finite number \(N=N(x,n)\), computes the next omitted
term, proves a rational ratio bound
\[
  \rho=\frac{|x|}{N+1}<1,
\]
and outputs
\[
  \left[
    S_N(x)-\frac{|x|^N/N!}{1-\rho},
    S_N(x)+\frac{|x|^N/N!}{1-\rho}
  \right].
\]
\end{definition}

\begin{definition}[Complex power-series exponential]\label{def:exp-power-series-function}
\lean{ComputableAnalysis.exp.ps}
\uses{def:function-raw,def:complex-raw,def:exp-power-series-real}
\leanok
The complex power-series representation is defined on all rational complex
inputs and uses the same power-series enclosure machinery in complex boxes.
Its real-axis restriction is
\[
  \code{ComputableAnalysis.exp.realPowerSeries}.
\]
\end{definition}

\begin{theorem}[Power-series validity package]\label{thm:exp-power-series-valid}
\lean{ComputableAnalysis.ExpProofs.PowerSeriesValid}
\uses{def:exp-power-series-real}
For every rational input \(x\), the intervals in
\(\code{ComputableAnalysis.expPowerSeries}\) are ordered, nested after
refinement, and have widths tending to zero.  The formal proof is separated
into finite
obligations:
\[
  \rho_n<1,\qquad
  \text{tail widths shrink to }0,\qquad
  \text{successive displayed boxes are nested}.
\]
This separation keeps the executable definition readable while making the
remaining proof work local to rational inequalities.
\end{theorem}

\section{Euler-Limit Exponential}

\begin{definition}[Euler-limit exponential]\label{def:exp-euler-real}
\lean{ComputableAnalysis.expEuler}
\uses{def:raw-real}
\leanok
For \(x\in\Q\), define
\[
  \lim_{m\to\infty}\left(1+\frac{x}{m}\right)^m
\]
as the Euler-limit exponential raw real by evaluating the finite product at
a stage-dependent integer \(m=m(n)\), and displaying it with an explicit
rational radius.
\end{definition}

\begin{definition}[Euler-limit function representation]\label{def:exp-euler-function}
\lean{ComputableAnalysis.exp.euler}
\uses{def:function-raw,def:exp-euler-real}
\leanok
The complex Euler-limit representation is the entire function
\[
  z\longmapsto \lim_{m\to\infty}\left(1+\frac{z}{m}\right)^m.
\]
Its finite stage computes the rational complex power
\[
  \left(1+\frac{z}{m}\right)^m
\]
and wraps it in a complex error box.  The real-axis restriction is
\[
  \code{ComputableAnalysis.exp.realEuler}.
\]
\end{definition}

\begin{theorem}[Euler-limit validity package]\label{thm:exp-euler-valid}
\lean{ComputableAnalysis.ExpProofs.EulerValid}
\uses{def:exp-euler-real}
For every rational \(x\), the Euler-limit interval algorithm represents a
raw real.  The current proof boundary reduces validity to the finite
nesting estimates for
\[
  \left(1+\frac{x}{m(n)}\right)^{m(n)}
\]
together with the explicit stage-radius shrinkage.
\end{theorem}

\section{Logarithm as an Integral}

\begin{definition}[Positive reciprocal kernel]\label{def:positive-reciprocal-kernel}
\lean{ComputableAnalysis.RatFun.oneOverXOnPositiveInterval}
\uses{def:function-on-interval}
\leanok
On a certified positive rational interval \(0<a\le x\le b\), the rational
function
\[
  x\longmapsto \frac1x
\]
is a valid function on the interval.  The positivity certificate is the
domain-apartness data needed for interval division.
\end{definition}

\begin{definition}[Logarithmic integral]\label{def:log-integral}
\uses{def:positive-reciprocal-kernel,def:monotone-integral}
For \(x>0\), define logarithm constructively by
\[
  \operatorname{Log}(x)=\int_1^x \frac{\dd t}{t}.
\]
On a finite positive interval this is built by integrating the reciprocal
kernel.  For \(x<1\), the orientation convention is
\[
  \operatorname{Log}(x)=-\int_x^1 \frac{\dd t}{t}.
\]
The formal target is a raw real-valued function whose domain is the positive
rational line, with the FTC proof identifying its derivative as \(1/x\).
\end{definition}

\begin{definition}[Exponential from inverse logarithmic integral]
\label{def:exp-from-log-integral}
\lean{ComputableAnalysis.exp.LogIntegralInverseBranch}
\uses{def:log-integral,def:function-on-interval}
\leanok
A finite inverse branch packages a positive source interval, the logarithmic
integral on that interval, and an inverse raw function.  If \(y\) lies in the
certified range, the branch returns the unique positive \(x\) with
\[
  \int_1^x \frac{\dd t}{t}=y.
\]
This is the integral definition of exponential on that branch.  The
associated raw function is
\[
  \code{ComputableAnalysis.exp.fromLogIntegral}.
\]
\end{definition}

\begin{definition}[Integral-inverse exponential on an interval]
\label{def:exp-from-log-integral-interval}
\lean{ComputableAnalysis.exp.FromLogIntegralOnInterval}
\uses{def:exp-from-log-integral,def:function-on-interval}
\leanok
For a rational interval of \(y\)-values, an inverse-log-integral branch gives
a function on that interval:
\[
  y\longmapsto \operatorname{Log}^{-1}(y).
\]
The branch data contains exactly the domain proof, pointwise validity proof,
and inverse certificate needed to treat this as a normal function on an
interval.
\end{definition}

\section{Agreement of the Definitions}

\begin{definition}[Three-way agreement target]\label{def:exp-three-way-agreement}
\lean{ComputableAnalysis.exp.representationsAgree}
\uses{def:exp-power-series-function,def:exp-euler-function,def:exp-from-log-integral}
\leanok
The main comparison statement says that the power-series function, the
Euler-limit function, and the inverse-log-integral function agree on every
common certified domain.
In Lean this is deliberately phrased as a package of pairwise agreement
statements, because the three proofs use different estimates.
\end{definition}

\begin{theorem}[Agreement from the three estimates]
\label{thm:exp-three-way-agreement-from-estimates}
\lean{ComputableAnalysis.exp.representationsAgree_of_estimates}
\uses{def:exp-three-way-agreement}
\leanok
If the power-series/Euler estimate, the power-series/integral-inverse
estimate, and the Euler/integral-inverse estimate have all been proved, then
the three representations agree.
\end{theorem}

\begin{theorem}[Power series agrees with Euler limit]
\label{thm:exp-power-series-euler}
\lean{ComputableAnalysis.ExpProofs.expPowerSeries_eq_expEuler_of_centersOverlap}
\uses{def:exp-power-series-real,def:exp-euler-real}
For each rational \(x\), the power-series and Euler-limit raw real
algorithms are equivalent once their displayed intervals are valid and the
center sequences overlap at every requested precision.  For \(x=1\), this
gives the comparison between the two algorithms for \(e\), recorded as
\[
  \code{ComputableAnalysis.ExpProofs.EPowerSeriesEqEuler}.
\]
\end{theorem}

\begin{theorem}[Power series agrees with inverse logarithmic integral]
\label{thm:exp-power-series-log-integral}
\lean{ComputableAnalysis.exp.powerSeries_equiv_logIntegralInverse_on_interval}
\uses{def:exp-power-series-function,def:exp-from-log-integral-interval,thm:convex-ftc}
\leanok
On a rational interval, the comparison is reduced to uniqueness for the
initial-value problem
\[
  f'=f,\qquad f(0)=1.
\]
The power-series exponential solves this equation by termwise
differentiation and the inverse-log-integral exponential solves it by the
inverse-function theorem applied to
\[
  \operatorname{Log}'(x)=1/x.
\]
The current Lean theorem states the final reduction from these two solution
certificates to function equivalence.
\end{theorem}

\begin{definition}[Euler product versus logarithmic integral estimate]
\label{def:exp-euler-log-integral-estimate}
\lean{ComputableAnalysis.exp.EulerProductIntegralEstimate}
\uses{def:exp-euler-function,def:exp-from-log-integral}
\leanok
The direct comparison between
\[
  \left(1+\frac{x}{m}\right)^m
\]
and the inverse logarithmic integral is a finite product estimate.  It is the
constructive version of comparing a multiplicative step
\[
  u\longmapsto u\left(1+\frac{x}{m}\right)
\]
with the accumulated integral of \(1/t\).  This estimate supplies the
Euler/integral-inverse component of the three-way agreement package.
\end{definition}

\section{Formalization Order}

\begin{enumerate}
\item Complete the real raw validity proofs for
\(\code{ComputableAnalysis.expPowerSeries}\) and
\(\code{ComputableAnalysis.expEuler}\).  The intended closing theorem is
\[
  \code{ComputableAnalysis.ExpProofs.complete_of_remainders}.
\]

\item Promote the reciprocal kernel into the logarithm function
\[
  \operatorname{Log}(x)=\int_1^x \frac{\dd t}{t}
\]
on positive rational inputs, using \cref{def:monotone-integral} and
\cref{thm:convex-ftc}.

\item Prove that the logarithmic integral is strictly increasing on positive
intervals, construct finite inverse branches, and assemble them into the
interval representation of exponential obtained from the logarithmic
integral.

\item Prove \(\operatorname{Log}'(x)=1/x\) and the inverse-branch derivative identity, then
close \cref{thm:exp-power-series-log-integral} through the uniqueness theorem
for \(f'=f,\ f(0)=1\).

\item Prove the Euler product estimate, close the Euler/log-integral
comparison, and combine the three pairwise estimates with
\cref{thm:exp-three-way-agreement-from-estimates}.

\item Define the final logarithm as the integral function and prove the
inverse laws on certified domains:
the logarithm of the exponential represents the original input, and the
exponential of the logarithm represents the original positive input.
\end{enumerate}
